% Source: https://artofproblemsolving.com/community/c3h1783549_math_problem
% Stable identity: aops-topic-1783549-q01
Classify the $26$ dice that are visible by how many of their faces lie on the surface:
\begin{itemize}
\item the $8$ corner dice show $3$ faces, which meet at a vertex of the die and are pairwise adjacent;
\item the $12$ edge dice (middle of each edge) show $2$ adjacent faces;
\item the $6$ face-centre dice show $1$ face;
\end{itemize}
and the centre die is hidden.

\begin{asy}
usepackage("amsmath");
size(4.5cm);
defaultpen(fontsize(10pt));
pair v=(0.45,0.32);
pen[] col={palegreen,paleblue,pink};
int ext(int i){return (i==0||i==2)?1:0;}
for(int i=0;i<3;++i) for(int j=0;j<3;++j){
 filldraw((i,j)--(i+1,j)--(i+1,j+1)--(i,j+1)--cycle,col[ext(i)+ext(j)]);
 filldraw(((i,3)+j*v)--((i+1,3)+j*v)--((i+1,3)+(j+1)*v)--((i,3)+(j+1)*v)--cycle,col[ext(i)+ext(j)]);
 filldraw(((3,i)+j*v)--((3,i)+(j+1)*v)--((3,i+1)+(j+1)*v)--((3,i+1)+j*v)--cycle,col[ext(i)+ext(j)]);
}
draw((0,0)--(3,0)--(3,3)--(0,3)--cycle,linewidth(0.9));
draw((0,3)--(0,3)+3v--(3,3)+3v--(3,0)+3v--(3,0),linewidth(0.9));draw((3,3)--(3,3)+3v,linewidth(0.9));
filldraw(box((0,-1.0),(0.4,-0.6)),pink);label("$3$ faces",(0.5,-0.8),E);
filldraw(box((0,-1.8),(0.4,-1.4)),paleblue);label("$2$ faces",(0.5,-1.6),E);
filldraw(box((0,-2.6),(0.4,-2.2)),palegreen);label("$1$ face",(0.5,-2.4),E);
label("\mbox{}",(min(currentpicture,user=true)+max(currentpicture,user=true))/2,black); // keeps the text after the figure black when LaTeX typesets the labels inline
\end{asy}

Three faces of a die meeting at a vertex contain no pair of opposite faces, so they show three different numbers, with sum at least $1+2+3=6$; and $1,2,3$ can indeed meet at a vertex, since no two of them are opposite ($1+6=2+5=3+4=7$). Two adjacent faces show different numbers with sum at least $1+2=3$, which is attainable. A single face shows at least $1$. Each die can be turned independently, so all these minima are achieved together, and the smallest total is
\[
8\cdot6+12\cdot3+6\cdot1=48+36+6=\boxed{90}.
\]
